% 从旧数学附录迁出的多项式函数说明；不新增算法。
\section{多项式带余除法：先反转再求逆}\label{knowledge-poly-division}
\texttt{PolynomialDivision::divide(a,b)} 返回规范化的商、余式，零多项式用空向量。
去掉输入高次零后，除数必须非零；低次被除式直接返回零商和原被除式。
设 $A=\deg a$、$B=\deg b$、$k=A-B+1$，反转后余式从 $x^k$ 起，故
\[
 \operatorname{rev}(q)=\operatorname{rev}(a)\,
 \operatorname{rev}(b)^{-1}\pmod{x^k}.
\]
除数最高项非零，反转后可求逆。取前 $k$ 项再反转得到商，最后用 $r=a-qb$ 得余式。
P4512 要求定长输出，使用程序负责补零；Library Checker 输出去高次零后的长度，零项对应空行。
输入各长度不超过 $2^{22}$ 是当前完整卷积容量的充分界，时间 $O(N\log(N+1))$、空间 $O(N)$。
独立长除法、五十万项构造与官方本地 checker 已双模式验证；线上待补。
\section{级数开根：前导零和自由尾项}\label{knowledge-poly-sqrt}
\texttt{FpsSqrt::sqrt(a,n)} 仅使用前 $n$ 项，不足补零。全零返回全零；
首个非零项次数 $k$ 为奇数，或其系数非二次剩余，则返回无解。
否则设 $a=x^kh$，从较小的非零模平方根开始，迭代
\[
 b_{\rm new}=\frac{b+h/b}{2},\qquad
 b_{\rm new}^2-h=\frac{(b^2-h)^2}{4b^2}.
\]
误差最低次数翻倍，故求出 $h$ 的根前 $n-k$ 项后左移 $k/2$ 即可。
所得根末尾还有 $k/2$ 个不受平方同余约束的系数，本接口统一取零。
因此一般截断问题不应说成“只有两个解”。返回 optional；$n=0$ 返回有值的空向量。
$n\le 2^{22}$；时间 $O(n\log(n+1))$ 加一次标量模开根，额外空间 $O(n)$。
P5205 的 $a_0=1$ 返回首项 $1$，一般前导零与无解另由 Library Checker 使用程序展示。
二次平方参考、最低次数边界、五十万项闭式及官方本地 checker 已双模式验证；线上待补。
\section{ln 与 exp：只处理所需前缀}\label{knowledge-poly-ln-exp}
对数要求 $a_0=1$，由 $(\log a)'=a'/a$ 求导、求逆、积分得到；
指数要求 $a_0=0$，从 $b=1$ 开始使用 $b\gets b(1+a-\log b)$ 倍增。
仅求前 $n$ 项时，先把对数输入截到 $n$ 项，避免输入的无关尾项造成变换容量超限。
积分需 $1,\ldots,|a|$ 都可逆；固定素数域中要求 $|a|<998244353$。
当前 $n\le 2^{22}$ 是变换容量充分界，log/exp 为 $O(n\log(n+1))$ 时间、$O(n)$ 空间。
测试使用独立三角系数递推，并包含超过 NTT 容量的长输入只求一项、非二次幂截断与五十万项闭式。
\section{幂函数的大指数：保留三份指数信息}\label{knowledge-poly-power}
\texttt{FpsPower::power(a,exponent,n)} 读入十进制数字串，计算模 $x^n$ 的非负整数次幂。
$n=0$ 返回空；零次幂统一返回单位多项式，包含 $0^0$。
若 $k>0$，设 $a=x^sc h$、$h(0)=1$、$c\ne0$；全零或 $sk\ge n$ 时结果全零，否则
\[
 a^k=x^{sk}c^k\exp(k\log h)\pmod{x^n}.
\]
同时保留 $k\bmod p$、$k\bmod(p-1)$ 和 $\min(k,n)$：第一份用于对数系数，第二份用于 $c^k$，
第三份用于真实零次幂及位移。当前 $n\le2^{22}<p$，故 $h(x)^p=h(x^p)\equiv1\pmod{x^n}$，
单位首项部分只依赖 $k\bmod p$。指数余数为零并不表示原指数为零。
判定 $sk\ge n$ 先用 $k>(n-1)/s$ 比较，通过后才乘，避免指数乘法溢出。
指数位数 $L$，时间 $O(L+n\log(n+1))$、额外空间 $O(n)$。
P5245 的常数项为 1，P5273 与 Library Checker 允许一般首项；共用主程序必须将指数读作字符串。
独立截断乘法、十万位指数、五十万项二项式及官方本地 checker 已双模式验证；线上待补。
\section{集合变换：逆变换对模数的要求}\label{knowledge-fwt-inverse}
OR 变换是子集和，单层为 $(u,v)\mapsto(u,u+v)$；AND 是超集和，单层为 $(u+v,v)$。
它们只需加减，模数可以为偶数或合数。XOR 变换使用字符
$(-1)^{\operatorname{popcount}(s\mathbin{\&}i)}$；单层 $(u+v,u-v)$ 连做两次等于乘 2，
故当前接口要求奇数模数，以 $(p+1)/2$ 为逆 2，不依赖费马小定理。
长度为正的二次幂，长度 1 可用；三种卷积不能与不相交子集卷积混淆。
独立字符定义式、下标对枚举、合数模数及最大规模完整驱动已双模式核验。
\section{浮点 FFT：舍入前先约束最坏误差}\label{knowledge-float-fft}
复数变换采用正指数，逆变换使用共轭根并除以变换长度。
长度为一时没有蝶形操作；位逆序从一开始，避免负移位。
单位根按各自角度直接计算，避免长链旋转因子连乘。
本节要求二进制至少 53 位有效精度、最近舍入、禁用 fast-math，
且单位根复数绝对误差不超过 $8u$，其中 $u=2^{-53}$。
三角函数精度是平台前提，不由 C++ 类型名称保证；本机已逐个核验最大网格的全部半圆根。

设变换长度 $N=2^\ell$，两输入欧氏范数乘积为 $S$。
\href{https://www.daemonology.net/papers/fft.pdf}{Percival 定理 5.1}
对三次 radix-2 FFT 卷积给出误差界
\[
 S\left((1+u)^{3\ell}(1+\sqrt5u)^{3\ell+1}
 (1+8u)^{3\ell}-1\right).
\]
整数入口限制 $N\le2^{22}$、
$\max|a_i|\max|b_i|\sqrt{|a||b|}\le2\cdot10^{12}$，
从而每项误差小于 $0.166$；使用 llroundl 同时处理正负系数。
超过此限制不能只因浮点结果接近整数就宣称正确。

任意模卷积先把有符号输入归一化，再拆成三个 11 位块。
结果长度最多 $2^{19}$。
每路 $S\le2047^2(2^{18}+1)$，误差小于 $0.079$。
六次正变换复用，九次逆变换分别舍入后用整数模合并。
此处以较大的常数换取明确的 53 位精度边界，不假定 long double 更宽。
模数可为一或合数；空输入返回空，常数乘法合法。
百万次数整数卷积、完整 int 输入、最大长度拆系数卷积已用闭式与精确整数参考核验；
本地验证不等于在线 AC。
